This report presents the design and implementation of \textbf{IceVentes}, a web application for managing ice-block sales, developed during an internship at the National Fisheries Office (Office National des Pêches, ONP). In the fishing sector, block ice is essential for preserving fish; its sale to professionals was previously tracked manually, without centralisation or traceability. The project aims to digitise this process through a modern, secure and easy-to-use web application.

The solution is built on the Python \textbf{Django} framework, a responsive interface using Bootstrap~5 and HTMX, a relational database (SQLite in development, PostgreSQL in production) and office-document exports (Excel, PDF). It covers the whole business cycle: authentication and user management with three roles (Administrator, Agent, Cashier); configuration of an ice-factory reference (code, city, unit price in MAD/Kg); management of requesters (customers, natural or legal persons, buyers or sellers); sale entry with automatic generation of a unique code \texttt{BG\_<ref>/<year> <number>} and automatic computation of quantity or total price; and finally a multi-criteria sales history with exports and a dashboard summarising activity (number of sales, quantity sold, revenue, active requesters) together with a monthly trend chart.

Particular care was taken regarding data consistency (unique codes, a single active reference, reliable numbering under concurrency), access security (role-based control, brute-force protection) and user-experience quality. The report describes the project context, the requirements specification, the system design and its actual implementation, illustrated by the produced interfaces.
